\section{Preprocessing Pipeline} The preprocessing code is intentionally lightweight. The script \verb|scripts/preprocess_data.sh| calls \verb|src.data.midi_preprocess|, which loads raw MIDI files, optionally remaps drum pitches, trims empty tracks, filters small or broken files, computes deduplication signatures, and writes cleaned MIDI files and optional manifests.



The main design choice is to avoid transformations that MidiTok already handles. \verb|src/data/midi_preprocess.py| explicitly avoids velocity binning, note/time quantization, and tempo discretization during cleaning. Those choices are deferred to the tokenizer configuration, which keeps the boundary between raw MIDI sanitation and symbolic tokenization clearer.

The current preprocessing defaults include:
\begin{itemize}
	\item coarse GM drum remapping via \verb|drum_map = gm_coarse|;
	\item minimum note count of 50 note onsets;
	\item minimum length of 4 estimated bars;
	\item deduplication by a PPQ-independent note signature;
	\item signature sensitivity to non-drum program numbers and note durations;
	\item no PPQ rebasing by default.
\end{itemize}

For the current xLSTM experiments, split creation reuses part of this preprocessing logic for analysis rather than writing a separate cleaned MIDI corpus. It parses candidate MIDI files, rejects parse errors and empty files, records top instrument labels and note counts, remaps drum pitches and trims empty tracks before computing the normalized content signature, removes duplicate content within and across music IDs, and then performs leakage-safe train/validation/test assignment. The command wrapper is:
\begin{quote}
	\small\ttfamily
	scripts/create\_dataset\_splits.sh \textless subset\_size\textgreater{} [genre\_csv] [out\_dir]\\
	\hspace*{2em}[seed] [jobs] [genres\_csv] [log\_level] [quiet] [extra\_args...]
\end{quote}
